\section*{अध्याय \ref{ch:g1:landscapes-around-us} --- हमारे आसपास के भूदृश्य}

\begin{solution}{exo:g1:landscapes-around-us:1}
\omterm{def:g1:landscapes-around-us:landform}{पर्वत}, \omterm{def:g1:landscapes-around-us:landform}{पहाड़ी}, \omterm{def:g1:landscapes-around-us:landform}{घाटी}, \omterm{def:g1:landscapes-around-us:landform}{मैदान}।
\end{solution}

\begin{solution}{exo:g1:landscapes-around-us:2}
\omterm{def:g1:landscapes-around-us:landform}{पर्वत}: वह \omterm{def:g1:landscapes-around-us:landform}{पहाड़ी} से ज़्यादा ऊँचा और ज़्यादा खड़ा होता है।
\end{solution}

\begin{solution}{exo:g1:landscapes-around-us:3}
प्राकृतिक: \omterm{def:g1:landscapes-around-us:water}{नदी}, जंगल, रेतीला किनारा। मानव-निर्मित: सड़क, पुल,
घर।
\end{solution}

\begin{solution}{exo:g1:landscapes-around-us:4}
जहाँ ज़मीन समुद्र से मिलती है।
\end{solution}

\begin{solution}{exo:g1:landscapes-around-us:5}
उदाहरण के लिए: प्राकृतिक --- \omterm{def:g1:landscapes-around-us:landform}{पर्वत}, ढलानों पर का जंगल या नाला;
मानव-निर्मित --- कोई घर या गाँव का गिरजाघर।
\end{solution}

\begin{solution}{exo:g1:landscapes-around-us:6}
\omterm{def:g1:landscapes-around-us:landform}{पर्वत}, जिसकी चोटी पर बर्फ़ है। \omterm{def:g1:landscapes-around-us:water}{नदी} \omterm{def:g1:landscapes-around-us:landform}{घाटी} की तली में, \omterm{def:g1:landscapes-around-us:landform}{पर्वतों} और \omterm{def:g1:landscapes-around-us:landform}{पहाड़ी}
के बीच बहती है।
\end{solution}

\begin{solution}{exo:g1:landscapes-around-us:7}
\omterm{def:g1:landscapes-around-us:water}{नदी} ऊँची ज़मीन (\omterm{def:g1:landscapes-around-us:landform}{पहाड़ियों} और \omterm{def:g1:landscapes-around-us:landform}{पर्वतों}) से नीचे किसी \omterm{def:g1:landscapes-around-us:water}{झील} या समुद्र की ओर
बहती है। पानी हमेशा ढलान की ओर नीचे बहता है, इसलिए वह किसी \omterm{def:g1:landscapes-around-us:landform}{पहाड़ी} की
चोटी तक ऊपर नहीं बह सकता।
\end{solution}

\begin{solution}{exo:g1:landscapes-around-us:8}
पेड़ हमारे पास है, अग्रभूमि में; \omterm{def:g1:landscapes-around-us:landform}{पहाड़ी} बहुत दूर है, पृष्ठभूमि में। दूर
की चीज़ें पास की चीज़ों से छोटी दिखती हैं।
\end{solution}

\begin{solution}{exo:g1:landscapes-around-us:9}
अग्रभूमि: चित्र के सबसे नीचे ऊँची घास और कुछ झाड़ियाँ।
मध्यभूमि: \omterm{def:g1:landscapes-around-us:landform}{मैदान} में बल खाती \omterm{def:g1:landscapes-around-us:water}{नदी} और उसके किनारे-किनारे पेड़, खेत, एक
फ़ार्म और एक सड़क।
पृष्ठभूमि: क्षितिज तक फैली समतल ज़मीन, और आसमान।
\end{solution}

\begin{solution}{exo:g1:landscapes-around-us:10}
क्योंकि उसे लोगों ने बनाया है: उन्होंने ज़मीन साफ़ की, जुताई की और गेहूँ
बोया। लोगों के बिना वहाँ गेहूँ का खेत नहीं, बल्कि घास, झाड़ियाँ या जंगल
होता।
\end{solution}

\begin{solution}{exo:g1:landscapes-around-us:11}
नहीं बदले: \omterm{def:g1:landscapes-around-us:landform}{पहाड़ी} और \omterm{def:g1:landscapes-around-us:water}{नदी}। लोगों ने जोड़े: खेत, ज़्यादा घर, सड़क और पुल
(और ज़्यादातर जंगल काट दिया गया है)।
\end{solution}
